\documentclass[../../main.tex]{subfiles}
\graphicspath{{../../image/}{../../figure/}}
\begin{document}

\section{能量估计}\label{section:ch15-能量估计}

此处估计\eqref{align:ch15-15421}右端最后一项，它源于我们按\eqref{align:ch15-15418}固定速度 $s_i$。目标是证明
\begin{align}
\int_{\bar t}^{T}\int_{-\infty}^{\infty}\bigl(|\partial_xw_i|+|\partial_x\omega_i|\bigr)|\omega_i+s_iw_i|\,dx\,dt=O(\delta_0^2).
\label{align:ch15-1571}
\end{align}
该证明依赖能量估计，技术上较繁，且不如前两节讨论那样富有启发性。因而这里只给出梗概。读者可在\cref{section:ch15-注记}所引文献中找到细节。

由于可作变量替换 $x\mapsto x-\lambda_i^*t$，可不失一般性地假设 $\lambda_i^*=0$，从而 $\sigma_i=O(\delta_0)$。

除由\eqref{align:ch15-15419}定义的 $\theta$ 外，我们还采用“截断”函数
\begin{align}
\eta(r)=
\begin{cases}
0&|r|\leqslant\frac35\delta_1,\\
1&|r|\geqslant\frac45\delta_1,
\end{cases}
\qquad|\eta'(r)|\leqslant\frac{20}{\delta_1},\qquad|\eta''(r)|\leqslant\frac{100}{\delta_1^2},
\label{align:ch15-1572}
\end{align}
以及 $\bar\eta(r)=\eta(|r|-\frac15\delta_1)$。记 $\eta_i=\eta(\omega_i/w_i)$，$\bar\eta_i=\bar\eta(\omega_i/w_i)$。并取 $\delta_0\ll\delta_1\ll1$。

\begin{lemma}\label{lemma:ch15-能量估计预备}
当 $|\omega_i/w_i|\geqslant\frac35\delta_1$ 时，
\begin{align}
\begin{cases}
|w_i|\leqslant\dfrac5{2\delta_1}|\partial_xw_i|+O(\delta_0)\displaystyle\sum_{j\neq i}|w_j|,\\[2mm]
|\omega_i|\leqslant2|\partial_xw_i|+O(\delta_0)\displaystyle\sum_{j\neq i}|w_j|,
\end{cases}
\label{align:ch15-1573}
\end{align}
而当 $|\omega_i/w_i|\leqslant\delta_1$ 时，
\begin{align}
|\partial_xw_i|\leqslant2\delta_1|w_i|+O(\delta_0)\sum_{j\neq i}|w_j|.
\label{align:ch15-1574}
\end{align}
\end{lemma}

\begin{proof}[证明梗概]
把\eqref{align:ch15-15418}中的 $\partial_xU$ 与 $\partial_tU$ 代入\eqref{align:ch15-1521}，得
\begin{align}
\sum_{j=1}^n\omega_jS_j+\sum_{j=1}^nw_jAS_j=\sum_{j=1}^n\partial_xw_jS_j+\sum_{j=1}^nw_j\partial_xS_j.
\label{align:ch15-1575}
\end{align}
左端乘以 $S_i^\top$，回顾 $S_i^\top S_i=1$，$S_i^\top\partial_xS_i=0$，并利用\eqref{align:ch15-1549}，得
\begin{align}
\omega_i-\sigma_iw_i-\partial_xw_i&=\sum_{j\neq i}\left\{[\partial_xw_j-\omega_j]S_i^\top S_j+w_jS_i^\top\partial_xS_j-w_jS_i^\top AS_j\right\}\nonumber\\
&=O(\delta_0)\sum_{j\neq i}\bigl(|w_j|+|\partial_xw_j-\omega_j|\bigr).
\label{align:ch15-1576}
\end{align}
断言\eqref{align:ch15-1573}与\eqref{align:ch15-1574}由对上式的细致分析得到。证明完成。
\end{proof}

由于 $|\omega_i+s_iw_i|$ 当 $|\omega_i/w_i|\leqslant\delta_1$ 时为零、否则以 $|\omega_i|$ 为界，故
\begin{align}
|\omega_i+s_iw_i|\leqslant|\bar\eta_i\omega_i|\leqslant\bar\eta_i\left[2|\partial_xw_i|+O(\delta_0)\sum_{j\neq i}|w_j|\right].
\label{align:ch15-1577}
\end{align}

因此，
\begin{align}
(|\partial_xw_i|+|\partial_x\omega_i|)|\omega_i+s_iw_i|&\leqslant2\bar\eta_i|\partial_xw_i|^2+2\bar\eta_i|\partial_xw_i|\,|\partial_x\omega_i|+\sum_{j\neq i}\bigl(|w_j\partial_xw_i|+|w_j\partial_x\omega_i|\bigr)\nonumber\\
&\leqslant3\eta_i|\partial_xw_i|^2+\bar\eta_i|\partial_x\omega_i|^2+\sum_{j\neq i}\bigl(|w_j\partial_xw_i|+|w_j\partial_x\omega_i|\bigr).
\label{align:ch15-1578}
\end{align}

如\cref{section:ch15-横波相互作用}所示，\eqref{align:ch15-1578}右端第三项在 $(-\infty,\infty)\times(\bar t,T)$ 上的积分为 $O(\delta_0^2)$。于是为验证\eqref{align:ch15-1571}，只需证明
\begin{align}
\int_{\bar t}^{T}\int_{-\infty}^{\infty}\eta_i|\partial_xw_i|^2\,dx\,dt=O(\delta_0^2),
\label{align:ch15-1579}
\end{align}
\begin{align}
\int_{\bar t}^{T}\int_{-\infty}^{\infty}\bar\eta_i|\partial_x\omega_i|^2\,dx\,dt=O(\delta_0^2).
\label{align:ch15-15710}
\end{align}

建立\eqref{align:ch15-1579}的第一步是用 $2\eta_iw_i$ 乘\eqref{align:ch15-1526}$_1$，在 $(-\infty,\infty)$ 上积分，并分部积分。这给出
\begin{align}
\int_{-\infty}^{\infty}\left\{\partial_t(\eta_iw_i^2)+\partial_x(\eta_i\sigma_i)w_i^2-(\partial_t\eta_i+2\sigma_i\partial_x\eta_i-\partial_x^2\eta_i)w_i^2+2\eta_i(\partial_xw_i)^2+4(\partial_x\eta_i)w_i\partial_xw_i\right\}\,dx=\int_{-\infty}^{\infty}2\eta_iw_i\varphi_i\,dx.
\label{align:ch15-15711}
\end{align}
因此，
\begin{align}
2\int_{-\infty}^{\infty}\eta_i|\partial_xw_i|^2\,dx=-\frac{\mathrm{d}}{\mathrm{d}t}\int_{-\infty}^{\infty}\eta_iw_i^2\,dx+\int_{-\infty}^{\infty}(\partial_t\eta_i+2\sigma_i\partial_x\eta_i-\partial_x^2\eta_i)w_i^2\,dx\nonumber\\
+2\int_{-\infty}^{\infty}\eta_i\sigma_iw_i\partial_xw_i\,dx-4\int_{-\infty}^{\infty}(\partial_x\eta_i)w_i\partial_xw_i\,dx+2\int_{-\infty}^{\infty}\eta_iw_i\varphi_i\,dx.
\label{align:ch15-15712}
\end{align}

我们进而估计上式右端。回顾 $\eta_i$ 的定义并利用\eqref{align:ch15-1526}，简单计算后得
\begin{align}
(\partial_t\eta_i+2\sigma_i\partial_x\eta_i-\partial_x^2\eta_i)w_i^2=\eta_i'(\psi_iw_i-\varphi_i\omega_i)+2\eta_i'w_i(\partial_xw_i)\partial_x(\omega_i/w_i)-\eta_i''w_i^2[\partial_x(\omega_i/w_i)]^2.
\label{align:ch15-15713}
\end{align}

进而，利用\eqref{align:ch15-1573}$_1$ 且因 $\sigma_i=O(\delta_0)\ll\delta_1$，
\begin{align}
\left|2\int_{-\infty}^{\infty}\eta_i\sigma_iw_i\partial_xw_i\,dx\right|\leqslant\int_{-\infty}^{\infty}\eta_i|\partial_xw_i|^2\,dx+O(\delta_0)\int_{-\infty}^{\infty}\sum_{j\neq i}|w_j\partial_xw_i|\,dx.
\label{align:ch15-15714}
\end{align}

在 $\eta_i'\neq0$ 的范围上，$|\omega_i/w_i|<\delta_1$，从而\eqref{align:ch15-1574}适用。于是
\begin{align}
|(\partial_x\eta_i)w_i\partial_xw_i|&=|\eta_i'w_i(\partial_xw_i)\partial_x(\omega_i/w_i)|\nonumber\\
&\leqslant O(1)|w_i\partial_x\omega_i-\omega_i\partial_xw_i|+O(\delta_0)\sum_{j\neq i}\bigl(|w_j\partial_xw_i|+|w_j\partial_x\omega_i|\bigr).
\label{align:ch15-15715}
\end{align}

现联合\eqref{align:ch15-15712}、\eqref{align:ch15-15713}、\eqref{align:ch15-15714}、\eqref{align:ch15-15715}，并在 $(\bar t,T)$ 上积分所得不等式。这给出形如
\begin{align}
\int_{\bar t}^{T}\int_{-\infty}^{\infty}\eta_i|\partial_xw_i|^2\,dx\,dt&\leqslant\int_{-\infty}^{\infty}(\eta_iw_i^2)(x,\bar t)\,dx\nonumber\\
&\quad+O(1)\int_{\bar t}^{T}\int_{-\infty}^{\infty}\bigl(|w_i\psi_i|+|w_i\varphi_i|+|\omega_i\varphi_i|\bigr)\,dx\,dt\nonumber\\
&\quad+O(1)\int_{\bar t}^{T}\int_{-\infty}^{\infty}|w_i\partial_x\omega_i-\omega_i\partial_xw_i|\,dx\,dt\nonumber\\
&\quad+O(\delta_0)\int_{\bar t}^{T}\int_{-\infty}^{\infty}\sum_{j\neq i}\bigl(|w_j\partial_xw_i|+|w_j\partial_x\omega_i|\bigr)\,dx\,dt\nonumber\\
&\quad+O(1)\int_{\bar t}\int_{|\omega_i/w_i|<\delta_1}^{T}|w_i\partial_x(\omega_i/w_i)|^2\,dx\,dt
\label{align:ch15-15716}
\end{align}
的估计。

由\eqref{align:ch15-15420}、\eqref{align:ch15-1527}、\eqref{align:ch15-1551}、\eqref{align:ch15-1561}与\eqref{align:ch15-1562}，得\eqref{align:ch15-15716}右端为 $O(\delta_0^2)$，从而验证\eqref{align:ch15-1579}。

估计\eqref{align:ch15-15710}由类似程序建立。细节见\cref{section:ch15-注记}所引文献。

\end{document}
